\section{Proofs of the uniform certificate tools}
\label{p14:comp:certificate-tools}

All definitions and parameter conditions are those of
Section~\ref{p14:cert:section}. The counts are geometric, and all
constants are uniform in the size of the family being summed.

\subsection{Proof of Lemma~\ref{p1:core:rankin}}
\label{p14:comp:certificate-rankin}

\begin{proof}
Write uniquely $v=k s a^2$, where $k=K_Y(v)$ and $s$ is the
squarefree product of the primes at most $Y$ occurring to odd order.
The inequality $1\le (C_1N/(ksa^2))^\sigma$ on the range being
counted gives an upper bound
\[
 (C_1N)^\sigma
 \left(\sum_{k\le T}k^{-\sigma}\right)
 \left(\sum_{a\ge1}a^{-2\sigma}\right)
 \prod_{\ell\le Y}(1+\ell^{-\sigma}).
\]
Enlarging the sum over $k$ to all positive integers only increases it.
The three factors are respectively $O_\sigma(T^{1-\sigma})$,
$\zeta(2\sigma)$, and at most
$\exp(\sum_{j\le Y}j^{-\sigma})$. This proves the bound.
\end{proof}

\subsection{Proof of Lemma~\ref{p1:core:moderate}}
\label{p14:comp:certificate-moderate}

\begin{proof}
There are at most $rNB^2$ choices for the anchor and its two offsets.
The product $Q$ has at most $m$ prime factors, and each has at most
$rB$ destinations. Once these assignments are fixed, the Chinese
remainder theorem leaves at most $1+N/D_j$ possibilities for centre
$j$. At most two of the $D_j$ exceed $N$, since $Q\le9N^2$.
For those two modules, replacing $1+N/D_j$ by $2N/D_j$ costs at most
$(6B)^2$. Therefore
\[
 \prod_{j\ne\mathrm{anchor}}(1+N/D_j)
 \le 2^{r-1}(6B)^2N^{r-1}T^{-2}.
\]
This proves \eqref{p1:core:moderate-count}; allowing collisions and
incompatible assignments only enlarges the count.

After multiplication by $(up)^r/r!$, the remaining series is, up to
a fixed factor,
\[
 N^{-2\beta}B^{m+4}\sum_{r\ge0}\frac{z^r r^{m+1}}{r!},
 \qquad z=2uNp=O_u(1).
\]
Since $r^{m+1}\le(m+1)!e^r$, this series is at most
$N^{-2\beta}B^{m+4}(m+1)!\exp(ez)$. Finally
\[
 m=\left(\frac2{A_0}+o(1)\right)\frac{H}{\log H},\qquad
 \log(B^m(m+1)!)=\left(\frac4{A_0}+o(1)\right)H.
\]
This yields \eqref{p1:core:moderate-baseline}. The factorial here
divides the orderings of a set of distinct windows; it is not an
assumption on the growth of a fixed-order arithmetic constant.
\end{proof}

\subsection{Proof of Lemma~\ref{p1:core:degrees}}
\label{p14:comp:certificate-degrees}

\begin{proof}
For an outgoing edge, choose two anchor offsets in $B^2$ ways. Each of
the at most $m$ primes is either unused or assigned to one of $B$
offsets in the receiver. A qualifying code has modulus $D>N$ and
therefore determines at most one receiving centre in $I_N$. This gives
$B^2(B+1)^m=N^{2/A_0+o(1)}$ possibilities.

For the incoming degree, fix the receiving window $y$ and let
$\Omega_y$ be the primes greater than $Y$ dividing any of its values.
Writing $R=\lceil B\log(3N)/\log Y\rceil$, we have $|\Omega_y|\le R$.
A code chooses two anchor offsets, at most $m$ primes from $\Omega_y$,
and one of the two offsets for each chosen prime. For a qualifying
code their product exceeds $N$, so its congruences determine at most
one anchor centre. Thus the number of possible incoming neighbours is
at most
\[
 B^2\sum_{j=0}^m\binom Rj2^j
 \le B^2(m+1)(2eR/m)^m
 =N^{2/A_0+o(1)}.
\]
The displayed binomial bound holds for large $N$, since $m\le R/2$.
Nonqualifying codes may be included in the number of codes; the
one-centre conclusion is applied only to qualifying ones. No receiving
offset needs to be chosen in this incoming count: the receiver has
already been fixed. Symmetrisation adds the two degree bounds.
\end{proof}

\subsection{Proof of Lemma~\ref{p1:core:generator}}
\label{p14:comp:certificate-generator}

\begin{proof}
There are at most $n(4\Delta)^{h-1}$ connected vertex sets of size
$h\ge2$: encode a canonical rooted spanning tree by its plane-tree
shape, its root and the choices of neighbours along its edges. The
number of plane-tree shapes is at most $4^{h-1}$. Allowing repeated
vertices only enlarges this count.

Decompose each admissible set into its induced connected components.
Singleton components must be marked. Dropping disjointness and the
absence of edges between the other components bounds the generating
function by
$\exp(Mz+\sum_{h\ge2}n(4\Delta)^{h-1}z^h)$.
In the exponential, the orderings of distinct components cancel the
factorial denominator; repeated components merely contribute additional
nonnegative terms. Summing the geometric series proves the claim.
\end{proof}
\section{From fixed orders to all primitive families}
\label{p1:core:section}

The estimates at each fixed order do not control the growing order needed
to reconstruct a distribution. We now obtain that uniformity without
estimating the growth of the fixed-order constants. The argument has two
parts. A large squarefree kernel produces congruences which count the
possible centres. When too much of this kernel is received by one other
window, it produces an edge in a graph of small degree. A generating
function then controls all large families at once. Only a finite set of
small cardinalities is left to the fixed-order estimates.

Throughout this section, a family of distinct windows is \emph{primitive}
if its spaces are not in direct sum but every proper subfamily is in
direct sum. Connectedness refers to the canonical decomposition into direct-sum
components in Lemma~\ref{p1:alg:components}.
All families in this section consist of windows in the regular set $G_N$.
The quantities $q(I)$, $c(I)$, $\rho(I)$ and $\rho_*(I)$ retain their
definitions from Section~\ref{p1:alg:section}. All weights and all auxiliary
parameters in this section are fixed before $N$ tends to infinity.

\subsection{The centred estimate at each fixed order}

The small-cardinality part of the argument needs an estimate for absolute
centred coefficients, including incompatible affine systems. Positive
probabilities alone do not give this estimate: a joint event can be
impossible while its centred coefficient is nonzero. For separated windows
we remove this difficulty by using a homogeneous majorant. Contacts require
a different argument, since their homogeneous probabilities can be large.

For $x\in G_N$, let $H_x$ be the indicator of the homogeneous event that
every character in $S_x$ takes the value $1$. Equivalently, the random
multiplicative function is constant on the $L+1$ values in window $x$. Set
\begin{equation}
 W(I)=\E\prod_{x\in I}H_x
      =p^{|I|}2^{\rho(I)}.
 \label{p1:core:homogeneous-weight}
\end{equation}
All homogeneous systems are compatible. Moreover, the full-support Fourier
formula of Lemma~\ref{p1:alg:fourier}, after taking absolute values, gives
\begin{equation}
 |c(I)|\le W(I).
 \label{p1:core:homogeneous-majorant}
\end{equation}
Here the geometric good-pair class includes pairs whose prescribed affine
run-start events are incompatible; denote this enlarged class by
$\mathcal A_{\rm raw}$.

\begin{lemma}[Absolute centred mass at a fixed order]
\label{p1:core:fixed-centred}
For every fixed integer $r\ge3$ and $\varepsilon>0$,
\begin{equation}
 \sum_{\substack{C\subseteq G_N,\ |C|=r\\C\ \mathrm{connected}}}
       |c(C)|\ll_{r,\varepsilon}N^{-4/5+\varepsilon}.
 \label{p1:core:fixed-centred-bound}
\end{equation}
The sum over all $r$-element families containing a contact, with no
connectedness restriction, is $O_{r,\varepsilon}(N^{-1+\varepsilon})$.
\end{lemma}
\begin{proof}
First suppose that all windows are separated. Call a family residual if
its relation space is neither zero nor exactly the relation space of one
geometric good pair, extended by zero on the other blocks. A connected
family of at least three windows is residual. We prove the stronger bound
obtained by summing $W(C)$ over all residual separated families of size $r$.

Proposition~\ref{p1:fixed:propagation} applies to the homogeneous targets
as well as to the prescribed affine targets. Its rank increment
$\dim(S_x\cap\sum_{y\in I}S_y)$ is geometric, and the conditional
probability is $p$ times the corresponding power of $2$ in either case.
Consequently a seed may be completed to any fixed total cardinality at
a cost $N^{\varepsilon}$, after reducing the error exponent used for
the seed. The following four cases exhaust residual families.

If a family contains a dependent pair outside $\mathcal A_{\rm raw}$,
the first seed estimate of Proposition~\ref{p1:arith:fixed-homogeneous-seeds}
gives total weight $O_\varepsilon(N^{-4/5+\varepsilon})$. Propagation
adds the other $r-2$ windows at the stated cost. If it contains an
original primitive subfamily of size $h\ge3$, with $h\le r$,
Proposition~\ref{p1:fixed:primitive} gives the same bound for its
homogeneous weight. There are only finitely many such values of $h$,
so propagation again gives the required estimate.

Next suppose that the family contains two distinct geometric good pairs.
Their union has three or four windows. By
Proposition~\ref{p1:arith:fixed-homogeneous-seeds}, the total homogeneous
weight of these unions is
\[
 O_\varepsilon(N^{-5/4+\varepsilon}+L^2N^{-4/3})
       =O_\varepsilon(N^{-5/4+\varepsilon}).
\]
Completing the union to $r$ windows costs only $N^{\varepsilon}$.
This case includes a pair sharing one endpoint with the other pair.

The remaining family has a unique geometric good pair $e$ and an
additional relation. Compress $e$ exactly, and choose a minimal dependent
subfamily of the resulting blocks. If it avoids $e$, it is one of the
original primitive seeds already considered. Otherwise it consists of
the compressed pair and $h-1$ original windows. For $h=2$, the
pair-and-target seed estimate of
Proposition~\ref{p1:arith:fixed-homogeneous-seeds} gives
$O_\varepsilon(N^{-4/3+\varepsilon})$ for channel $2:3$ and
$O_\varepsilon(N^{-5/4+\varepsilon})$ for channel $3:4$.
For $h\ge3$, Proposition~\ref{p1:fixed:compressed} gives
\[
 O_{h,\varepsilon}(N^{-\kappa_b+1/5+\varepsilon}),
 \qquad \kappa_b=2-1/b,\quad b\in\{3,4\}.
\]
The largest of these bounds is $O(N^{-22/15+\varepsilon})$.
Propagation is applied after replacing the compressed seed by its original
windows. The number of seeds and placements in a family of fixed size is
bounded in terms of $r$. These cases prove the separated estimate by
\eqref{p1:core:homogeneous-majorant}. No compatibility restriction has
been introduced in this argument.

For contacts, return to the actual affine events. Expanding the centred
product gives
\begin{equation}
 |c(C)|\le\sum_{U\subseteq C}p^{r-|U|}q(U).
 \label{p1:core:contact-expansion}
\end{equation}
Fix $h=|U|$ and write $u=r-h$. If $U$ already contains a contact,
the positive contact seed followed by
Proposition~\ref{p1:fixed:propagation} gives
\[
 \sum_{\substack{|U|=h\\U\ \mathrm{contains\ a\ contact}}}q(U)
       \ll_{h,\varepsilon}N^{-1+\varepsilon}.
\]
The remaining $u$ centres and their factors $p^u$ cost at most
$(Np)^u=O_r(1)$. If $U$ is separated but $C$ contains a contact,
at least one endpoint of a contact is among the $u$ new centres.
For fixed $U$ there are at most $O_r(BN^{u-1})$ such completions:
either a new centre is within $O(B)$ of $U$, or two new centres are
within $O(B)$ of each other. The cases where no completion is possible
contribute zero. Since
\[
 \sum_{|U|=h}q(U)
       =\E\binom{Z_N^\circ}{h}=O_h(1)
\]
by the fixed moments in Theorem~\ref{p1:fixed:moments}, these
completions contribute $O_r(B/N)$. Summing over the finitely many
values of $h$ proves the contact assertion. The moments used here are
proved from the positive affine classification before this lemma; they
do not depend on any all-cardinality centred estimate.
\end{proof}

\subsection{Primitive families at all cardinalities}

\begin{theorem}[Primitive families at all cardinalities]
\label{p1:core:primitive}
For every fixed $w>0$ and $\varepsilon>0$,
\begin{equation}
 \sum_{\substack{C\subseteq G_N\ \mathrm{primitive}\\|C|\ge3}}
 w^{|C|}\bigl(q(C)+|c(C)|\bigr)
 \ll_{w,\varepsilon}N^{-4/5+\varepsilon}.
 \label{p1:core:primitive-bound}
\end{equation}
The sum includes overlapping windows and incompatible affine systems.
\end{theorem}

\subsection{One rank loss for a primitive family}

If $C$ is primitive, the projection of its relation space onto any one
block is injective. Indeed, an element in the kernel of that projection
would be a relation supported on a proper subfamily. Since a regular
window has rank $L$, a primitive family of $r\ge3$ windows satisfies
\begin{equation}
 1\le\rho(C)\le L,\qquad
 q(C)\le p^{r-1},\qquad c(C)=q(C)-p^r.
 \label{p1:core:primitive-rank}
\end{equation}
The last identity follows by expanding the centred product: every proper
subfamily has joint probability equal to the product of its marginals.
Consequently $|c(C)|\le p^{r-1}$, including when $q(C)=0$.

A nonzero relation of a primitive family uses every block. Its projection
on a regular window corresponds to a nonempty even selection of the
$B=L+1$ values in that window. It therefore selects at least two distinct
values there. This elementary fact is what permits a two-value counting
argument, rather than a one-value friability estimate.

The kernel, moderate-certificate, saturation graph and generating-function
estimates used next are stated in Section~\ref{p14:cert:section} and
proved in Section~\ref{p14:comp:certificate-tools}. Their $p$-weighted
forms are used here on the regular set.

\subsection{The finite-small-order, uniform-large-order split}

Take
\begin{equation}
 A_0=40,\quad\beta=\frac{24}{25},\quad\sigma=\frac{49}{50},
 \quad\delta=\frac1{1250},\quad\theta=\frac2{A_0}=\frac1{20}.
 \label{p1:core:primitive-parameters}
\end{equation}
These values leave a strict margin in all subsequent inequalities.
For a fixed $w>0$, substitute $z=wpN^{\delta/2}$ into
Lemma~\ref{p1:core:generator}. The three quantities in its exponent
or denominator satisfy
\[
 Mz\le N^{-\delta/2+o(1)},\qquad
 \Delta z\le N^{-1+\theta+\delta/2+o(1)},\qquad
 n\Delta z^2\le N^{-1+\theta+\delta+o(1)}.
\]
They tend to zero. Consequently, for every fixed integer $s\ge3$,
\begin{equation}
 \sum_{\substack{C\ \mathrm{primitive},\ |C|\ge s\\
                         C\ \mathrm{without\ a\ moderate\ certificate}}}
 w^{|C|}\bigl(q(C)+|c(C)|\bigr)
 \ll_w N^{1-s\delta/2}.
 \label{p1:core:primitive-tail}
\end{equation}
Indeed, removing the tilt contributes $N^{-s\delta/2}$ and the
primitive rank bound contributes at most $2p^{-1}=O(N)$.
For the moderate families, Lemma~\ref{p1:core:moderate} and the same
rank bound give $N^{1-2\beta+4/A_0+o(1)}=N^{-41/50+o(1)}$.

\begin{proof}[Proof of Theorem~\ref{p1:core:primitive}]
Use $s=5000$, so the nonmoderate tail in
\eqref{p1:core:primitive-tail} is $O_w(N^{-1})$.
For each of the finitely many sizes $3\le r<s$, the fixed-order
primitive estimate in Proposition~\ref{p1:fixed:primitive} gives
$O_{r,\varepsilon}(N^{-4/5+\varepsilon})$
for the positive mass of separated families. Compatible contacts and
their fixed-order extensions contribute $O_{r,\varepsilon}(N^{-1+\varepsilon})$.
Lemma~\ref{p1:core:fixed-centred} controls $|c(C)|$,
including contacts and incompatibility. Alternatively, for a primitive
one can use $|c(C)|\le q(C)+p^r$ and the fixed-order support count.
Replacing its fixed weight by $w^r$ costs only a constant on this
finite set of sizes. Add these estimates to the two large-size bounds.
No fixed-order estimate has been applied with $r=r(N)$.
\end{proof}

\section{One distinguished pair after compression}
\label{p1:core:compression-section}

The approximation will keep the contribution of one good pair.
To control additional relations above it, the pair must first be
treated as a single affine block. Centring this new block is different
from centring its two original indicators; we retain the distinction
until the final identity of this section.

For an affine-compatible good pair $e=\{x,y\}$ of channel
$(a,b)=(2,3)$ or $(3,4)$,
let $J_e=J_xJ_y$, $q_e=\E J_e$ and $V_e=S_x+S_y$.
Its exact compression has centre $X_e=bx$, offset
$h_e=ay-bx=O(B)$, and values
\[
 U_e=b\mathcal V_x\cup a\mathcal V_y,
 \qquad \mathcal V_x=\{x-1,x,\ldots,x+L-1\}.
\]
Thus $U_e$ is a set of integer values, while $V_e$ is a vector space.
Lemma~\ref{p1:arith:compression} gives an
injective even-selection representation on the distinct values of
$U_e$, width $O(B)$, at most $2B$ values, and
\begin{equation}
 r_e=\dim V_e\ge\kappa_bL-O(1),\qquad
 q_e=2^{-r_e}\ll N^{-\kappa_b},\qquad \kappa_b=2-1/b.
 \label{p1:core:compressed-rank}
\end{equation}
There are $B^{O(1)}$ possible motifs; for a fixed motif its centre
determines the pair subject to the integral congruences. The affine
event itself is preserved, without an additional random sign from the
multipliers.

For a set $A$ of original windows disjoint from $e$, put
\[
 q(e,A)=\E\left[J_e\prod_{z\in A}J_z\right],\qquad
 \widehat c(e,A)=\E\left[(J_e-q_e)\prod_{z\in A}(J_z-p)\right].
\]
The configuration is primitive after compression when the spaces
$V_e,(S_z)_{z\in A}$ are primitive. Its number of compressed blocks
is $r=|A|+1$ and its number of original windows is $r+1$.

\begin{theorem}[A distinguished compressed pair at all cardinalities]
\label{p1:core:compressed}
For fixed $w>0$, let
\[
 \mathcal P_b(w)=\sum_{\substack{(e,A)\ \mathrm{primitive\ after\ compression}\\
                  e\ \mathrm{in\ channel}\ a:b,\ |A|\ge1}}
 w^{|A|+2}\{q(e,A)+|\widehat c(e,A)|+|c(e\cup A)|\}.
\]
Then, for each $\varepsilon>0$,
\begin{align}
 \mathcal P_3^{\mathrm{sep}}(w)&\ll_{w,\varepsilon}N^{-4/3+\varepsilon},&
 \mathcal P_4^{\mathrm{sep}}(w)&\ll_{w,\varepsilon}N^{-5/4+\varepsilon},
 \label{p1:core:compressed-separated}\\
 \mathcal P_b(w)&\ll_{w,\varepsilon}N^{-1+\varepsilon},&&b=3,4.
 \label{p1:core:compressed-all}
\end{align}
The superscript restricts all original windows to be separated. The
pair remains distinguished in each sum.
\end{theorem}

\begin{proof}
Projection onto any original singleton block is injective on the
primitive relation space. Thus its nullity is at most $L$, and
\begin{equation}
 \widehat c(e,A)=q(e,A)-q_ep^{r-1},\qquad
 q(e,A),\ |\widehat c(e,A)|\le q_ep^{r-2}
       \ll N^{2-\kappa_b}p^r.
 \label{p1:core:compressed-cost}
\end{equation}
For an incompatible configuration, the centred value is exactly
$-q_ep^{r-1}$. These estimates retain the saving from the pair rank.

Use the parameters \eqref{p1:core:primitive-parameters}. The two-anchor
certificate works on the compressed block as well: all its distinct
values lie in an interval of width $O(B)$ and are $O(N)$, so their
large kernels are coprime. Replace the moderate threshold by $C_2NB$
for a sufficiently large fixed constant, and put
$m_*=\lceil2\log(C_5N)/\log Y\rceil$, where $C_5$ bounds all
compressed values divided by $N$. With $P_N=B^{O(1)}$ possible
pair motifs, the ordered moderate host count is at most
\[
 P_N C_3^r r^2B^4N^rT^{-2}(C_4rB)^{m_*}.
\]
There is only one factor $P_N$, since only one compressed block is
present. Dividing by $r!$, using \eqref{p1:core:compressed-cost} and
summing the exponential series as before gives
\begin{equation}
 N^{2-\kappa_b-2\beta+4/A_0+o(1)}
       =N^{-\kappa_b+9/50+o(1)}
 \label{p1:core:compressed-moderate}
\end{equation}
for each of the first two masses. The factors $P_N$ and $C_4^{m_*}$
are $N^{o(1)}$.

For the nonmoderate sector use a typed graph whose vertices are all
original windows and all possible compressed good-pair blocks of the
fixed channel. It has $n_*=N^{1+o(1)}$ vertices,
$M_*=N^{1-\delta+o(1)}$ marked vertices, and maximum saturation degree
$\Delta_*=N^{\theta+o(1)}$. In the incoming and outgoing degree
counts a fixed motif adds only $B^{O(1)}$ choices. Every unmarked
vertex of a primitive configuration of at least three blocks has a
saturated neighbour. Lemma~\ref{p1:core:generator}, with
$z=wpN^{\delta/2}$, is still bounded because
$M_*z$, $\Delta_*z$, and $n_*\Delta_*z^2$ tend to zero with strict
power margins. We may enlarge the generating function to all typed
families, then retain \eqref{p1:core:compressed-cost} for the
configurations with exactly one distinguished pair. For each fixed
$s\ge3$, this bounds their tail by
\[
 O_w(N^{2-\kappa_b-s\delta/2}).
\]
At $s=5000$ this is $O_w(N^{-\kappa_b})$.

It remains to treat finitely many $2\le r<5000$. For separated
configurations with fixed $r\ge3$,
Proposition~\ref{p1:fixed:compressed} bounds their homogeneous mass by
$O_{r,\varepsilon}(N^{-\kappa_b+1/5+\varepsilon})$.
This bounds both $q$ and $|\widehat c|$, since the latter is at most
the homogeneous mass even when the full affine configuration is
incompatible. For $r=2$, the pair-and-target estimate of
Proposition~\ref{p1:arith:fixed-homogeneous-seeds} gives respectively
$O_\varepsilon(N^{-4/3+\varepsilon})$ and
$O_\varepsilon(N^{-5/4+\varepsilon})$ for the same homogeneous
majorant. Thus no phase restriction is needed in the separated estimate.

For the contact sector, Proposition~\ref{p1:fixed:propagation} bounds the
mass of the $r+1$ original windows by $O_{r,\varepsilon}(N^{-1+\varepsilon})$.
The number of possible distinguished pairs in such a fixed-size
family is bounded in terms of $r$. The deterministic host count for
these contacts is $N^{r-1+o(1)}$: the good pair costs $NB^{O(1)}$,
and a contact removes one remaining free centre. Multiplying by the
centering term $q_ep^{r-1}$ gives $N^{-\kappa_b+o(1)}$.
The identity in \eqref{p1:core:compressed-cost} therefore bounds
$|\widehat c|$ on contacts, including incompatible configurations.
Together with the two tails, these bounds prove
\eqref{p1:core:compressed-separated}--\eqref{p1:core:compressed-all}
for $q$ and $|\widehat c|$.

Finally, $A$ is a nonempty direct family, and expansion of
$(J_x-p)(J_y-p)$ gives the exact identity
\begin{equation}
 c(e\cup A)=\widehat c(e,A)
                -p\,c(\{x\}\cup A)-p\,c(\{y\}\cup A).
 \label{p1:core:centering-change}
\end{equation}
Each endpoint family is direct or primitive: a proper subset
containing the endpoint is contained in a proper direct family of
compressed blocks. A fixed original primitive family of size $r$
can be extended by a distinguished companion in at most $O(rB)$ ways,
since the deterministic good-pair degree is $O(B)$. Hence the sum of
the two corrections is bounded by a fixed multiple of
\[
 wpB\sum_{r\ge2}r w^r
             \sum_{\substack{D\ \mathrm{primitive}\\|D|=r}}|c(D)|.
\]
For $r\ge3$, absorb $r$ into a slightly larger fixed weight and use
Theorem~\ref{p1:core:primitive}; for $r=2$, use the total pair activity
$O(LN^{-2/3})$. The result is
$O_{w,\varepsilon}(pB(LN^{-2/3}+N^{-4/5+\varepsilon}))
 =O_{w,\varepsilon}(N^{-5/3+\varepsilon})$.
This is smaller than each claimed bound and proves the assertion for
the original centred coefficient as well.
\end{proof}

\subsection{The active structure outside a small event}

\begin{lemma}[A single pair is the only remaining active component]
\label{p1:core:exceptional}
There exists an event $E_*$ with
$\P(E_*)\ll_\varepsilon N^{-4/5+\varepsilon}$ such that, outside
$E_*$, all active spaces are a direct sum of singleton spaces and at
most one compressed good-pair space. In particular every active
subfamily has nullity at most $L/3+O(1)$.
\end{lemma}
\begin{proof}
Exclude an active contact, an active nongood dependent pair, an active
primitive family of at least three original windows, two distinct
active good pairs, and an active primitive family after compression of
a good pair. The pair and contact estimates, the two-pair seed bound in
Proposition~\ref{p1:arith:fixed-homogeneous-seeds}, and
Theorems~\ref{p1:core:primitive} and~\ref{p1:core:compressed}
bound the union by the asserted probability. All sums here are positive.

If no good pair is active, a minimal dependent active subfamily would
be one of the excluded families. If exactly one good pair is active,
compress it. A minimal dependent subfamily of the resulting blocks
either avoids the compressed block, giving an excluded original
primitive, or contains it, giving an excluded compressed primitive.
Thus the remaining compressed blocks are direct. The only rank loss
is that of the good pair, at most $L/3+O(1)$.
\end{proof}

\section{Connected coefficients with bounded compatible nullity}
\label{p1:core:nullity-section}

The preceding theorems concern minimal dependent families. The
boundary of the stopped expansion can contain larger connected
families. We first control all cardinalities when the compatible
nullity is less than two window ranks. We then remove the restriction
on that fixed nullity level over the growing range needed by the
expansion.

\subsection{A two-anchor estimate at every cardinality}

\begin{theorem}[All cardinalities below two ranks]
\label{p1:core:low-nullity}
For fixed $0\le\eta<2$, $w>0$ and $\varepsilon>0$,
\begin{equation}
 \sum_{\substack{C\subseteq G_N\ \mathrm{connected},\ |C|\ge3\\
                  \rho_*(C)\le\eta L}}
 w^{|C|}|c(C)|
 \ll_{\eta,w,\varepsilon}
 N^{-\min\{4/5,\,2-\eta\}+\varepsilon}.
 \label{p1:core:low-nullity-bound}
\end{equation}
\end{theorem}
\begin{proof}
Expansion of the centred product gives
\begin{equation}
 |c(C)|\le 2^{|C|}p^{|C|}2^{\rho_*(C)}
       \ll_\eta N^\eta(2p)^{|C|}.
 \label{p1:core:compatible-majorant}
\end{equation}
The Fourier expansion also shows that $c(C)\ne0$ implies a relation
meeting every block: centring removes the zero Fourier coordinate in
each factor. Such a relation again selects at least two values per
block. Therefore a family with nonzero coefficient and no moderate
certificate has only marked isolated vertices in the saturation graph.
No independence assertion about the edges of that graph is used.

With the parameters in \eqref{p1:core:parameter-conditions}, the
moderate sector, by Lemma~\ref{p1:core:moderate}, costs
$N^{\eta-2\beta+4/A_0+o(1)}$. The graph generator bounds its
nonmoderate tail beyond any fixed size $s$ by
$O_{\eta,w}(N^{\eta-s\delta/2})$.
For $0<\varepsilon<\min\{1,2-\eta\}$ choose
\[
 A_0>\max\{40,32/\varepsilon\},\qquad
 \beta=1-\varepsilon/16,\qquad
 \sigma=1-1/(2A_0).
\]
Then $\delta=\varepsilon/(32A_0)$, $2/A_0+\delta<1$, and
$2-2\beta+4/A_0<\varepsilon/4$. Choose a fixed integer $s$ with
$s\delta/2\ge\eta+1$. The moderate part has the required
$2-\eta$ exponent, the nonmoderate tail is $O(N^{-1})$, and
the finitely many connected sizes $3\le r<s$ are handled by
Lemma~\ref{p1:core:fixed-centred}.
Shrinking the margins there gives \eqref{p1:core:low-nullity-bound}.
Larger values of $\varepsilon$ follow from smaller ones.
\end{proof}

At $\eta=1$, this already gives the $4/5$ exponent at all sizes.
At $\eta=3/2$, however, the two-anchor estimate would give only
the exponent $1/2$. We need more anchors where a block carries many
independent relation coordinates.

\subsection{Projection dimensions and positive cores}

For a family $D$, let $\mathcal K(D)$ be the kernel of the sum map
$\bigoplus_{x\in D}S_x\to\sum_{x\in D}S_x$. The dimension of its
projection onto $S_x$ is
\begin{equation}
 d_x(D)=\rho(D)-\rho(D\setminus\{x\})
       =\dim\left(S_x\cap\sum_{y\in D\setminus\{x\}}S_y\right).
 \label{p1:core:projection-dimension}
\end{equation}
The kernel of that projection is exactly
$\mathcal K(D\setminus\{x\})$, which proves the identity.

If $t$ is even, a subspace of dimension at least $t-1$ in the
even-selection space contains a vector of weight at least $t$.
Indeed, choose $t-1$ pivot coordinates and set all of them equal to
one. The resulting vector has weight at least $t-1$, and its weight
is even. Call $D$ a \emph{$t$-core} when $d_x(D)\ge t-1$ for every
$x\in D$. At each of its blocks a relation therefore selects at least
$t$ distinct values. The selected relation may depend on the block.

\begin{proposition}[Positive mass of high-nullity cores]
\label{p1:core:positive-core}
Fix $\eta>1$, an even integer $t\ge\eta+2$, and $v>0$. Then
\begin{equation}
 \sum_{\substack{D\ \mathrm{compatible},\ D\ t\mathrm{-core}\\
                  L/2<\rho(D)\le\eta L}}
 v^{|D|}q(D)\ll_{\eta,t,v,\varepsilon}N^{-4/5+\varepsilon}
 \label{p1:core:positive-core-bound}
\end{equation}
for every $\varepsilon>0$, with no restriction on the cardinality.
\end{proposition}
\begin{proof}
We give the multiple-anchor count explicitly. Set
\[
\begin{gathered}
 A_0=8t,\qquad Y=(\log N)^{A_0},\qquad
 \beta=1-\frac1{4t},\qquad T=N^\beta,\\
 \sigma=1-\frac1{16t},\qquad\delta=\frac1{64t^2},\qquad
 m_t=\left\lceil\frac{t\log(3N)}{\log Y}\right\rceil.
\end{gathered}
\]
The number of marked windows remains $N^{1-\delta+o(1)}$.
At an unmarked anchor choose $t$ selected values. Their kernel
product $Q$ is squarefree, with $T^t<Q\le(3N)^t$ and at most $m_t$
prime factors. A moderate allocation has $D_j\le6NB$ for each
receiver. At most $t$ received modules exceed $N$, so its CRT cost is
\[
 \prod_{j\ne\mathrm{anchor}}(1+N/D_j)
 \le2^{r-1}(6B)^tN^{r-1}T^{-t}.
\]
The corresponding ordered host count is at most
$6^t rB^{2t}2^{r-1}N^rT^{-t}(rB)^{m_t}$.
Multiplication by $(up)^r/r!$ and the same exponential-series
estimate as before give, for fixed $u>0$,
\[
 \sum_r\frac{(up)^r}{r!}|\mathcal U_{r,t}|
 \le N^{-t\beta+2t/A_0+o(1)}.
\]
Here $\log(B^{m_t+2t}(m_t+1)!)=(2t/A_0+o(1))\log N$.
Since $q(D)\ll_\eta N^\eta p^{|D|}$, the positive moderate mass is
\begin{equation}
 N^{\eta-t\beta+2t/A_0+o(1)}
       =N^{\eta-t+1/2+o(1)}\le N^{-3/2+o(1)}.
 \label{p1:core:multi-anchor-gain}
\end{equation}

For the remaining cores, define saturation using $t$ anchor values.
The outgoing degree is bounded by $B^t(B+1)^{m_t}$. For the incoming
degree, the fixed receiver has at most
$R=\lceil B\log(3N)/\log Y\rceil$ possible primes; the codes are
bounded by
$B^t\sum_{j\le m_t}\binom Rj t^j$.
As in Lemma~\ref{p1:core:degrees}, a qualifying code has modulus
greater than $N$ and determines at most one centre. Both degrees are
$N^{t/A_0+o(1)}=N^{1/8+o(1)}$.
Every unmarked vertex of a nonmoderate $t$-core has a saturated
neighbour. Lemma~\ref{p1:core:generator}, tilted by $N^{\delta/2}$,
therefore bounds the positive mass beyond fixed size $s$ by
$O_{\eta,v}(N^{\eta-s\delta/2})$.
Choose $s\ge128t^2(\eta+1)$, so this is $O(N^{-1})$.

For a fixed size $h<s$, compatibility and $\rho(D)>L/2$ imply that
activation of $D$ occurs inside the event $E_*$ of
Lemma~\ref{p1:core:exceptional}. Indeed, its complement permits
nullity at most $L/3+O(1)$. For every fixed integer $k>1$, H\"older
and the bounded fixed moments of Theorem~\ref{p1:fixed:moments} give
\[
 \sum_{\substack{|D|=h,\ D\ \mathrm{compatible}\\\rho(D)>L/2}}q(D)
 \le\E\left[\binom{Z_N^\circ}{h};E_*\right]
 \le C_{h,k}\P(E_*)^{1-1/k}.
\]
Choose $k$ large enough that $4/(5k)<\varepsilon/3$, then use a
sufficiently small fixed margin in the bound for $\P(E_*)$.
The result is $O_{h,\varepsilon}(N^{-4/5+\varepsilon})$.
Only finitely many sizes $h$ are involved. Adding these to
\eqref{p1:core:multi-anchor-gain} and the saturated tail proves
\eqref{p1:core:positive-core-bound}. The distributional approximation
being proved in this paper is not used in this argument.
\end{proof}

\subsection{Removing blocks with small relation projections}

\begin{theorem}[Bounded compatible nullity on the growing range]
\label{p1:core:bounded-nullity}
For every fixed $\eta\ge0$, $w>0$, $\varepsilon>0$ and every
$K_N=o(L)$,
\begin{equation}
 \sum_{\substack{C\subseteq G_N\ \mathrm{connected}\\
          3\le|C|\le K_N,\ \rho_*(C)\le\eta L}}
 w^{|C|}|c(C)|\ll_{\eta,w,\varepsilon}N^{-4/5+\varepsilon}.
 \label{p1:core:bounded-nullity-bound}
\end{equation}
For $\eta>1$, choose a fixed even $t\ge\eta+2$ and set $d=t-2$.
The estimate holds uniformly whenever $K_N\le L/(2d)$. In particular,
for $\eta=3/2$ one may take $t=4$, $d=2$ and $K_N\le L/4$.
\end{theorem}
\begin{proof}
Theorem~\ref{p1:core:low-nullity} handles $\rho_*(C)\le L$.
For $L<\rho_*(C)\le\eta L$, choose a compatible subfamily
$A\subseteq C$ attaining $\rho(A)=\rho_*(C)$, using a fixed
lexicographic rule if needed. Remove from $A$ any block whose
relation projection has dimension at most $d$, until none remains.
Let $D$ be the resulting family. Equation~\eqref{p1:core:projection-dimension}
shows that
\[
 0\le\rho(A)-\rho(D)\le d(|A|-|D|)\le dK_N\le L/2.
\]
Thus $D$ is a nonempty compatible $t$-core with
$L/2<\rho(D)\le\eta L$.
If $r=|C|$ and $h=|D|$, the centred expansion and the rank bound give
\begin{equation}
 |c(C)|\le2^r p^r2^{\rho_*(C)}
        \le2^r(2^dp)^{r-h}q(D).
 \label{p1:core:peeling-weight}
\end{equation}
The elements of $C\setminus A$ are included in $p^{r-h}$; they need
not have been removed by the algorithm. This also covers incompatible
families $C$.

Sum \eqref{p1:core:peeling-weight} over all admissible cores and all
their supersets. Dropping restrictions on the latter only increases
the sum. With $n=|G_N|$ and $np\le Np=\lambda$, this gives
\begin{align*}
 \sum_{\substack{|C|\le K_N\\L<\rho_*(C)\le\eta L}}
 w^{|C|}|c(C)|
 &\le\sum_{D\ \mathrm{admissible}}(2w)^{|D|}q(D)
        (1+2^{d+1}wp)^{n-|D|}\\
 &\le e^{2^{d+1}w\lambda}
        \sum_{D\ \mathrm{admissible}}(2w)^{|D|}q(D).
\end{align*}
The last sum is bounded by Proposition~\ref{p1:core:positive-core}.
There is no factor for the deletion order: a deterministic choice
assigns one core to each family, and summing over all supersets already
overcounts the original family. The exponential factor is constant
because $d,w$ are fixed and $\lambda$ is bounded. Adding the sector
$\rho_*(C)\le L$ proves the theorem.
\end{proof}

For the required value $\eta=3/2$, the choices in the positive-core
proof are
$A_0=32$, $\beta=15/16$, $\sigma=63/64$, $\delta=1/1024$ and
$s=5120$. Its moderate mass is $N^{-2+o(1)}$, its saturated tail is
$O(N^{-1})$, and only fixed sizes below $5120$ use fixed moments.
This supplies the whole range containing
$K_N\asymp\log N/\log\log N$ without promoting a fixed-order
estimate to a growing-order one.

The same deletion argument gives the positive corollary
\[
 \sum_{\substack{D\ \mathrm{compatible},\ |D|\le K_N\\
                    L<\rho(D)\le\eta L}}
 w^{|D|}q(D)\ll_{\eta,w,\varepsilon}N^{-4/5+\varepsilon}.
\]
The centred form \eqref{p1:core:bounded-nullity-bound}, rather than
this additional corollary, is the estimate used by the stopped
expansion below.
